% GENERATED FILE. Edit canonical lessons, then run npm run generate:books.
% canonical-source-hash: fnv1a64:f85e58e5ac2ef19e
% canonical-lessons: LA-C499-delecto, LA-C499-delectatio, LA-C499-voluptas, LA-C499-oblecto, LA-C499-placeo

\chapter{Pleasure and delight}
\label{ch:la-pleasure-and-delight}
\hlchaptermodality{499}

\begin{chapteropening}
\textbf{By the end of this chapter:} \emph{I can talk about what pleases and delights me.}

Complete the last lesson of chapter 499: I can talk about what pleases and delights me.
\end{chapteropening}

\section[dēlectō, dēlectāre]{dēlectō, dēlectāre — to delight, charm}

\label{lesson:LA-C499-delecto}



\begin{quote}
Before the new one: say the Latin for to distrust, then the Latin for trusting.
\end{quote}

\subsection*{You'll want to know: dēlectō, dēlectāre}
\textbf{dēlectō, dēlectāre} — \textquotedblleft{}to delight, charm\textquotedblright{}.

\begin{quote}
Fābulae avī puerōs dēlectant.
\end{quote}

\begin{quote}
\emph{Grandfather's stories delight the children.}
\end{quote}

\emph{dēlectō, dēlectāre, dēlectāvī, dēlectātum}, with the accusative. English \emph{delight} comes from it through Old French.

\begin{cousinweb}
\textbf{dēlectō} is the frequentative of \emph{dēliciō} \textquotedblleft{}to allure, entice\textquotedblright{} (Lewis \& Short), built on \emph{dē-} and \emph{laciō} \textquotedblleft{}to entice\textquotedblright{}. English \emph{delight} comes from it through Old French.
\end{cousinweb}

\begin{grammarlens}[title={What you've built}]
One more word for telling, explaining and describing at B1.
\end{grammarlens}

\subsection*{Your turn}
\begin{itemize}
\raggedright
  \item \emph{Say it:} \emph{dēlectō, dēlectāre}
  \item \emph{Say it:} \emph{dēlectō, dēlectāre}, once more
  \item \emph{Say it:} a sentence of your own with \emph{dēlectō, dēlectāre}
  \item \emph{Recall:} say the Latin for to distrust, then the Latin for trusting, then say \emph{dēlectō, dēlectāre} again
\end{itemize}

\begin{tcolorbox}[breakable,colback=teal!4,colframe=teal!35!black,title={Before you move on}]
What does dēlectō, dēlectāre mean? (\textquotedblleft{}To delight, charm\textquotedblright{}.) Say it once more.
\end{tcolorbox}
\section[dēlectātiō, dēlectātiōnis]{dēlectātiō, dēlectātiōnis — delight, pleasure}

\label{lesson:LA-C499-delectatio}



\begin{quote}
Before the new one: say the Latin for trusting, then the Latin for to delight.
\end{quote}

\subsection*{You'll want to know: dēlectātiō, dēlectātiōnis}
\textbf{dēlectātiō, dēlectātiōnis} — \textquotedblleft{}delight, pleasure\textquotedblright{}.

\begin{quote}
Cantāre mihi magna dēlectātiō est.
\end{quote}

\begin{quote}
\emph{Singing is a great delight to me.}
\end{quote}

Feminine, third declension, from \emph{dēlectō}. English \emph{delectation}.

\begin{cousinweb}
\textbf{dēlectātiō} is the action noun of \emph{dēlectō} \textquotedblleft{}to delight\textquotedblright{} (Lewis \& Short). English \emph{delectation} comes from it.
\end{cousinweb}

\begin{grammarlens}[title={What you've built}]
One more word for telling, explaining and describing at B1.
\end{grammarlens}

\subsection*{Your turn}
\begin{itemize}
\raggedright
  \item \emph{Say it:} \emph{dēlectātiō, dēlectātiōnis}
  \item \emph{Say it:} \emph{dēlectātiō, dēlectātiōnis}, once more
  \item \emph{Say it:} a sentence of your own with \emph{dēlectātiō, dēlectātiōnis}
  \item \emph{Recall:} say the Latin for trusting, then the Latin for to delight, then say \emph{dēlectātiō, dēlectātiōnis} again
\end{itemize}

\begin{tcolorbox}[breakable,colback=teal!4,colframe=teal!35!black,title={Before you move on}]
What does dēlectātiō, dēlectātiōnis mean? (\textquotedblleft{}Delight, pleasure\textquotedblright{}.) Say it once more.
\end{tcolorbox}
\section[voluptās, voluptātis]{voluptās, voluptātis — pleasure, enjoyment}

\label{lesson:LA-C499-voluptas}



\begin{quote}
Before the new one: say the Latin for to delight, then the Latin for delight.
\end{quote}

\subsection*{You'll want to know: voluptās, voluptātis}
\textbf{voluptās, voluptātis} — \textquotedblleft{}pleasure, enjoyment\textquotedblright{}.

\begin{quote}
Sapiēns nōn voluptātem, sed virtūtem quaerit.
\end{quote}

\begin{quote}
\emph{The wise man seeks not pleasure but virtue.}
\end{quote}

Feminine, third declension. Pleasure of every kind, of the senses and of the mind. English \emph{voluptuous}.

\begin{cousinweb}
\textbf{voluptās} is built on the old adverb \emph{volup} \textquotedblleft{}pleasantly\textquotedblright{} (EtymDB), heard in the phrase \emph{volup est} \textquotedblleft{}it is a pleasure\textquotedblright{}; Lewis \& Short compares Greek \emph{elpo} \textquotedblleft{}to hope\textquotedblright{} instead. English \emph{voluptuous} comes from it.
\end{cousinweb}

\begin{grammarlens}[title={What you've built}]
One more word for telling, explaining and describing at B1.
\end{grammarlens}

\subsection*{Your turn}
\begin{itemize}
\raggedright
  \item \emph{Say it:} \emph{voluptās, voluptātis}
  \item \emph{Say it:} \emph{voluptās, voluptātis}, once more
  \item \emph{Say it:} a sentence of your own with \emph{voluptās, voluptātis}
  \item \emph{Recall:} say the Latin for to delight, then the Latin for delight, then say \emph{voluptās, voluptātis} again
\end{itemize}

\begin{tcolorbox}[breakable,colback=teal!4,colframe=teal!35!black,title={Before you move on}]
What does voluptās, voluptātis mean? (\textquotedblleft{}Pleasure, enjoyment\textquotedblright{}.) Say it once more.
\end{tcolorbox}
\section[oblectō, oblectāre]{oblectō, oblectāre — to amuse, entertain}

\label{lesson:LA-C499-oblecto}



\begin{quote}
Before the new one: say the Latin for delight, then the Latin for pleasure.
\end{quote}

\subsection*{You'll want to know: oblectō, oblectāre}
\textbf{oblectō, oblectāre} — \textquotedblleft{}to amuse, entertain\textquotedblright{}.

\begin{quote}
Senex sē librīs oblectat.
\end{quote}

\begin{quote}
\emph{The old man amuses himself with books.}
\end{quote}

\emph{oblectō, oblectāre, oblectāvī, oblectātum}. Often with \emph{sē}: to pass the time pleasantly.

\begin{cousinweb}
\textbf{oblectō} is built on \emph{ob-} and \emph{lactō} \textquotedblleft{}to allure\textquotedblright{}, a frequentative of \emph{laciō} \textquotedblleft{}to entice\textquotedblright{} (Lewis \& Short; EtymDB), the same root as \emph{dēlectō}.
\end{cousinweb}

\begin{grammarlens}[title={What you've built}]
One more word for telling, explaining and describing at B1.
\end{grammarlens}

\subsection*{Your turn}
\begin{itemize}
\raggedright
  \item \emph{Say it:} \emph{oblectō, oblectāre}
  \item \emph{Say it:} \emph{oblectō, oblectāre}, once more
  \item \emph{Say it:} a sentence of your own with \emph{oblectō, oblectāre}
  \item \emph{Recall:} say the Latin for delight, then the Latin for pleasure, then say \emph{oblectō, oblectāre} again
\end{itemize}

\begin{tcolorbox}[breakable,colback=teal!4,colframe=teal!35!black,title={Before you move on}]
What does oblectō, oblectāre mean? (\textquotedblleft{}To amuse, entertain\textquotedblright{}.) Say it once more.
\end{tcolorbox}
\section[placeō, placēre]{placeō, placēre — to please, be pleasing}

\label{lesson:LA-C499-placeo}



\begin{quote}
Before the new one: say the Latin for pleasure, then the Latin for to amuse.
\end{quote}

\subsection*{You'll want to know: placeō, placēre}
\textbf{placeō, placēre} — \textquotedblleft{}to please, be pleasing\textquotedblright{}.

\begin{quote}
Haec domus nova uxōrī meae placet.
\end{quote}

\begin{quote}
\emph{This new house pleases my wife.}
\end{quote}

\emph{placeō, placēre, placuī, placitum}, with the dative: \emph{mihi placet} — I like it; it seems good to me. English \emph{please} comes from it.

\begin{cousinweb}
\textbf{placeō} is a native Latin verb, a cousin of \emph{plācō} \textquotedblleft{}to calm\textquotedblright{} (EtymDB; Lewis \& Short). English \emph{please} and \emph{pleasure} come from it through Old French.
\end{cousinweb}

\begin{grammarlens}[title={What you've built}]
One more word for telling, explaining and describing at B1.
\end{grammarlens}

\subsection*{Your turn}
\begin{itemize}
\raggedright
  \item \emph{Say it:} \emph{placeō, placēre}
  \item \emph{Say it:} \emph{placeō, placēre}, once more
  \item \emph{Say it:} a sentence of your own with \emph{placeō, placēre}
  \item \emph{Recall:} say the Latin for pleasure, then the Latin for to amuse, then say \emph{placeō, placēre} again
\end{itemize}

\begin{tcolorbox}[breakable,colback=teal!4,colframe=teal!35!black,title={Before you move on}]
What does placeō, placēre mean? (\textquotedblleft{}To please, be pleasing\textquotedblright{}.) Say it once more.
\end{tcolorbox}
